\documentclass[solutions]{../cc}

\usepackage{tasks}
%\settasks{
%	label=\Alph*., % Labels will be A., B., C...
%	label-width=15pt,
%	column-sep=10pt,
%	after-item-skip=2pt
%}

\settasks{
	label=$\square$,       % Sets the default label to a box (or \Alph* $\square$)
	label-width=1.5em,     % Adds enough space for the box
	column-sep=20pt        % Space between columns
}

\usepackage{enumitem}

\begin{document}
	
	\makeheader{2025}{2026}
	\cctitle{1}{Bases de l'apprentissage automatique}
	\student
		
	\noindent \textbf{Consignes du QCM - Durée 20 minutes}:
	\begin{itemize}[nosep]
		\item Notes de cours et TD autorisés au format papier.
		\item Réponses multiples: Plusieurs réponses peuvent être correctes par question.
		\item Barème:
		\begin{itemize}[nosep]
			\item 2 points si toutes les bonnes réponses (et uniquement elles) sont cochées.
			\item 1 point si une partie des bonnes réponses est cochée, à condition qu'aucune mauvaise réponse ne soit sélectionnée.
			\item 0 point dès qu'une réponse fausse est cochée ou si la question est vide.
		\end{itemize}
		\item Notation: Le contrôle comporte 12 questions. Seuls vos 10 meilleurs scores seront retenus pour le calcul de la note finale sur 20.
	\end{itemize}

%   QUESTION 9
%	\noindent \begin{minipage}{\textwidth}
	%		\vspace{0.2cm}
	\begin{question}
		La méthode de bootstrapping:
		\begin{tasks}(1)
			\task Permet d'effectuer un estimation du risque vrai grâce aux données out-of-bag.
			\task Est une méthode de validation croisée.
			\task Divise les données d'apprentissage en plusieurs blocs disjoints.
			\task Ré-échantillonne les données d'apprentissage avec remise.
		\end{tasks}
	\end{question}
%	\end{minipage}

%   QUESTION 6
%	\noindent \begin{minipage}{\textwidth}
	%		\vspace{0.2cm}
	\begin{question}
		Pour un problème de régression $y = f(x) + \epsilon$, le coefficient de détermination $R^2$:
		\begin{tasks}(1)
			\task Mesure la fraction de la variance expliquée par le modèle.
			\task Est toujours supérieur à 1 si le modèle est bon.
			\task Vaut 1 si le coefficient de corrélation linéaire entre les prédictions $\hat{y}$ et les valeurs observées $y$ est égale à 1. 
			\task Vaut 0 si le modèle prédit la moyenne de $y$ pour tout $x$.
		\end{tasks}
	\end{question}
%	\end{minipage}

%   QUESTION 2
%	\noindent \begin{minipage}{\textwidth}
	%		\vspace{0.2cm}
	\begin{question}
		Cocher les affirmations vraies à propos des régression linéaires et logistiques.
		\begin{tasks}(1)
			\task La régression linéaire dispose d'une solution analytique.
			\task La régression logistique cherche à maximiser la log-vraisemblance.
			\task La régression logistique peut être résolue avec une méthode directe.
			\task La méthodes des moindres carrés cherche à maximiser la somme des écarts quadratiques.
		\end{tasks}
	\end{question}
%	\end{minipage}
	
%   QUESTION 8
%	\noindent \begin{minipage}{\textwidth}
	%		\vspace{0.2cm}
	\begin{question}
		La solution du problème de régression linéaire est donnée par $\hat{\beta} = (\tilde{X}^T \tilde{X})^{-1} \tilde{X}^T Y$ avec $\tilde{X} \in \mathcal{M}_{n,d} \in \mathbb{R}$ et $n > d$. Cette solution existe si:
		\begin{tasks}(1)
			\task Les lignes de $\tilde{X}$ sont linéairement indépendantes.
			\task La matrice $(\tilde{X}^T \tilde{X})$ est inversible.
			\task La première colonne de la matrice $\tilde{X}$ est un vecteur de $1$.
			\task La matrice $\tilde{X}$ est de plein rang.
		\end{tasks}
	\end{question}
%	\end{minipage}

%   QUESTION 3
%	\noindent \begin{minipage}{\textwidth}
	%		\vspace{0.2cm}
	\begin{question}
		La fonction logistique $\sigma$ possède les propriétés suivantes:
		\begin{tasks}(2)
			\task $\sigma(t) = 1 - \sigma(-t)$.
			\task Elle prend ses valeurs dans $]-1, 1[$.
			\task $\sigma(t) = -\sigma(-t)$.
			\task Elle est dérivable sur $\mathbb{R}$.
		\end{tasks}
	\end{question}
	%	\end{minipage}
	
%   QUESTION 12
%	\noindent \begin{minipage}{\textwidth}
	%		\vspace{0.2cm}
	\begin{question}
		Quelles sont les raisons d'une erreur d'approximation importante ?
		\begin{tasks}(1)
			\task Un mauvais choix d'espace d'hypothèses.
			\task L'absence de données de validation.
			\task Un manque de données d'apprentissage.
			\task Un bruit de mesure important.
		\end{tasks}
	\end{question}
%	\end{minipage}

%   QUESTION 11
%	\noindent \begin{minipage}{\textwidth}
	%		\vspace{0.2cm}
	\begin{question}
		En apprentissage automatique, lorsque l'on doit sélectionner un modèle ou un espace des hypothèses:
		\begin{tasks}(1)
			\task On choisit le modèle le plus complexe. Qui peut le plus, peut le moins.
			\task On choisit le modèle avec les meilleures métriques de performance sur les données d'apprentissage.
			\task On choisit un espace des hypothèses qui donne une explication plausible au problème avec une performance acceptable.
			\task On choisit un modèle avec une variance faible et un biais faible.
		\end{tasks}
	\end{question}
%	\end{minipage}

%   QUESTION 5
%	\noindent \begin{minipage}{\textwidth}
	%		\vspace{0.2cm}
	\begin{question}
		Que représente $\epsilon$ dans la relation $y = f(x) + \epsilon$ ?
		\begin{tasks}(2)
			\task Un biais d'approximation.
			\task Un bruit sur les caractéristiques observées.
			\task L'erreur de Bayes.
			\task Des variables explicatives non-observées.
		\end{tasks}
	\end{question}
%	\end{minipage}

%   QUESTION 1
%	\noindent \begin{minipage}{\textwidth}
%		\vspace{0.2cm}
		\begin{question}
			Dans la décomposition biais-variance de l'erreur quadratique d'un problème de régression, quels termes composent le risque vrai ponctuel ?
			\begin{tasks}(2)
				\task Le biais du modèle au carré.
				\task L'erreur de Bayes.
				\task La variance du modèle.
				\task L'écart de généralisation.
			\end{tasks}
		\end{question}
%	\end{minipage}

%   QUESTION 7
%	\noindent \begin{minipage}{\textwidth}
	%		\vspace{0.2cm}
	\begin{question}
		Le phénomène de sur-apprentissage est associé à:
		\begin{tasks}(1)
			\task Un écart de généralisation positif.
			\task Un manque relatif de données par rapport à la complexité du modèle.
			\task Un mauvais choix de métrique de performance.
			\task Un risque vrai qui devient négligeable par rapport au risque empirique.
		\end{tasks}
	\end{question}
%	\end{minipage}

%   QUESTION 4
%	\noindent \begin{minipage}{\textwidth}
%		\vspace{0.2cm}
		\begin{question}
			Qu'est-ce que le risque empirique ?
			\begin{tasks}(1)
				\task Une estimation de l'erreur de Bayes.
				\task La moyenne de la fonction de perte sur les données de test.
				\task La moyenne de la fonction de perte sur les données d'apprentissage.
				\task La performance d'une hypothèse sur l'ensemble complet des observations possibles.
			\end{tasks}
		\end{question}
%	\end{minipage}

%   QUESTION 10
%	\noindent \begin{minipage}{\textwidth}
	%		\vspace{0.2cm}
	\begin{question}
		En classification binaire, la courbe ROC:
		\begin{tasks}(1)
			\task Permet de visualiser la relation entre le taux de faux positifs et le taux de vrais positifs.
			\task Se calcule à partir des probabilités prédites, des classes vraies et du seuil de décision.
			\task Permet de visualiser l'influence du seuil de décision de l'algorithme de classification sur le compromis précision / rappel.
			\task Est révélatrice d'un bon apprentissage si la surface en dessous de la courbe tend vers 0.
		\end{tasks}
	\end{question}
	%	\end{minipage}
	
\end{document}